% Part: first-order-logic
% Chapter: sequent-calculus
% Section: structural-rules

\documentclass[../../../include/open-logic-section]{subfiles}

\begin{document}

\iftag{FOL}
      {\olfileid{fol}{seq}{srl}}
      {\olfileid{pl}{seq}{srl}}

\olsection{સંરચનાત્મક નિયમો}

સિક્વન્ટની ડાબી અને જમણી બાજુએ રહેલાં !!{sentence}sને ફરી ગોઠવવા દે
એવા થોડા નિયમો પણ જોઈએ. તાર્કિક નિયમો જે આધારવિધાનનાં !!{sentence}s પર
કાર્ય કરે છે તે છેક ડાબે અથવા છેક જમણે ઊભાં હોવાં જોઈએ; તેથી !!{sentence}sને
યોગ્ય સ્થાને ખસેડી શકાય એવો ``અદલાબદલી''નો નિયમ જોઈએ. ક્યારેક બે સમાન
!!{sentence}sને એકમાં ભેગાં કરી શકાય અને બંનેમાંથી ગમે તે બાજુએ
!!a{sentence} ઉમેરી શકાય એ પણ જરૂરી છે.

\subsection{શિથિલીકરણ}

\begin{defish}
\Axiom$ \Gamma \fCenter \Delta $
\RightLabel{\LeftR{\Weakening}}
\UnaryInf$ !A, \Gamma \fCenter \Delta$
\DisplayProof
\hfill
\Axiom$ \Gamma \fCenter \Delta$
\RightLabel{\RightR{\Weakening}}
\UnaryInf$ \Gamma \fCenter \Delta, !A$
\DisplayProof
\end{defish}

\subsection{સંકોચન}

\begin{defish}
\Axiom$ !A, !A, \Gamma \fCenter \Delta $
\RightLabel{\LeftR{\Contraction}}
\UnaryInf$ !A, \Gamma \fCenter \Delta$
\DisplayProof
\hfill
\Axiom$ \Gamma \fCenter \Delta, !A, !A$
\RightLabel{\RightR{\Contraction}}
\UnaryInf$ \Gamma \fCenter \Delta, !A$
\DisplayProof
\end{defish}

\subsection{અદલાબદલી}

\begin{defish}
\Axiom$ \Gamma, !A, !B, \Pi \fCenter \Delta $
\RightLabel{\LeftR{\Exchange}}
\UnaryInf$ \Gamma, !B, !A, \Pi \fCenter \Delta$
\DisplayProof
\hfill
\Axiom$ \Gamma \fCenter \Delta, !A, !B, \Lambda$
\RightLabel{\RightR{\Exchange}}
\UnaryInf$ \Gamma \fCenter \Delta, !B, !A, \Lambda$
\DisplayProof
\end{defish}

શિથિલીકરણ, સંકોચન અને અદલાબદલીનાં અનુમાનોની શ્રેણી ઘણી વાર બેવડી
અનુમાન-રેખાઓથી દર્શાવવામાં આવશે.

નીચે આપેલો ``કટ'' નામનો નિયમ કડક અર્થમાં જરૂરી નથી, પરંતુ તે
!!{derivation}sનો ફરી ઉપયોગ અને તેમનું સંયોજન ઘણું સહેલું બનાવે છે.
\begin{defish}
\[
\Axiom$ \Gamma \fCenter \Delta, !A$
\Axiom$ !A, \Pi \fCenter \Lambda $
\RightLabel{\Cut}
\BinaryInf$ \Gamma, \Pi \fCenter \Delta, \Lambda$
\DisplayProof
\]
\end{defish}

\end{document}
